%==============================================================================
% INSTRUÇÕES DE USO — Fichamentos Críticos
% Livro: Odontologia & Inteligência Artificial
% Conforme NBR 6023 (Referências) e NBR 10520 (Citações)
%==============================================================================

\chapter*{Instruções para Elaboração dos Fichamentos Críticos}
\addcontentsline{toc}{chapter}{Instruções — Fichamentos Críticos}

\section*{Propósito}
Os fichamentos críticos constituem o registro sistematizado da fundamentação
científica do livro. Cada artigo citado via \texttt{\textbackslash citacaoativa}
no texto deve possuir um fichamento correspondente, assegurando
rastreabilidade, auditabilidade e reprodutibilidade das fontes.

\section*{Estrutura do Arquivo}
Cada fichamento deve ser salvo como \texttt{fichamento-ODTXXX.tex}, onde
\texttt{XXX} é o identificador sequencial (001, 002, \ldots, 062).
\begin{itemize}
  \item Utilize o arquivo \texttt{template-fichamento.tex} como base.
  \item Substitua todos os campos entre colchetes pelos valores
        correspondentes do artigo fichado.
  \item A chave bibliográfica (\texttt{CHAVE\_BIB}) deve corresponder à
        entrada no arquivo \texttt{.bib} do projeto.
\end{itemize}

\section*{Seções Obrigatórias}
Cada fichamento deve conter obrigatoriamente as seguintes seções:

\begin{enumerate}
  \item \textbf{Identificação} — Dados bibliográficos completos conforme
        NBR 6023.
  \item \textbf{Objetivo} — Síntese do propósito do artigo em 3--5 linhas.
  \item \textbf{Metodologia} — Desenho do estudo, amostra, métricas e
        validação.
  \item \textbf{Resultados Principais} — Valores numéricos relevantes em
        lista itemizada.
  \item \textbf{Limitações} — Pelo menos duas limitações identificadas.
  \item \textbf{Contribuição para o Livro} — Relação explícita entre o
        artigo e o(s) capítulo(s) que ele fundamenta.
  \item \textbf{Análise Crítica} — Validade interna/externa,
        reprodutibilidade, relevância clínica, conflito de interesses e
        viés.
  \item \textbf{Citações Relevantes} — Transcrição de trechos citados no
        livro com indicação de página.
  \item \textbf{Mapeamento SDD} — Relação problema-entrada-saída-critérios
        conforme metodologia SDD do livro.
  \item \textbf{Referência ABNT} — Referência formatada conforme NBR 6023
        dentro de ambiente \texttt{verbatim}.
\end{enumerate}

\section*{Critérios de Qualidade}
\begin{itemize}
  \item \textbf{Rastreabilidade:} Toda afirmação no livro apoiada por
        \texttt{\textbackslash citacaoativa} deve ter fichamento
        correspondente.
  \item \textbf{Precisão:} Métricas numéricas (acurácia, sensibilidade,
        AUC, etc.) devem ser transcritas exatamente como publicadas.
  \item \textbf{Imparcialidade:} A análise crítica deve apontar tanto
        contribuições quanto limitações, sem viés de confirmação.
  \item \textbf{SDD:} O mapeamento SDD deve explicitar o problema
        odontológico, os tipos de dado de entrada/saída e os critérios de
        aceitação reportados no artigo.
  \item \textbf{Confiança:} Atribuir valor entre 0 (mínima) e 1 (máxima)
        refletindo a confiabilidade do estudo com base no desenho, tamanho
        amostral, validação e reprodutibilidade.
\end{itemize}

\section*{Verificação de DOIs}
Antes de finalizar o fichamento, verifique se o DOI do artigo está ativo
em \url{https://doi.org/}. DOIs que retornam \texttt{HTTP 404} ou
\texttt{HTTP 403} devem ser sinalizados no campo \textbf{DOI} com a
observação \texttt{(pendente -- possível DOI fabricado)} e reportados no
arquivo \texttt{00-doi-audit.tex}.

\section*{Integração ao Documento Principal}
Para incluir os fichamentos no livro, adicione no preâmbulo:
\begin{verbatim}
% Inclusão dos fichamentos críticos (apêndice documental)
%==============================================================================
% INSTRUÇÕES DE USO — Fichamentos Críticos
% Livro: Odontologia & Inteligência Artificial
% Conforme NBR 6023 (Referências) e NBR 10520 (Citações)
%==============================================================================

\chapter*{Instruções para Elaboração dos Fichamentos Críticos}
\addcontentsline{toc}{chapter}{Instruções — Fichamentos Críticos}

\section*{Propósito}
Os fichamentos críticos constituem o registro sistematizado da fundamentação
científica do livro. Cada artigo citado via \texttt{\textbackslash citacaoativa}
no texto deve possuir um fichamento correspondente, assegurando
rastreabilidade, auditabilidade e reprodutibilidade das fontes.

\section*{Estrutura do Arquivo}
Cada fichamento deve ser salvo como \texttt{fichamento-ODTXXX.tex}, onde
\texttt{XXX} é o identificador sequencial (001, 002, \ldots, 062).
\begin{itemize}
  \item Utilize o arquivo \texttt{template-fichamento.tex} como base.
  \item Substitua todos os campos entre colchetes pelos valores
        correspondentes do artigo fichado.
  \item A chave bibliográfica (\texttt{CHAVE\_BIB}) deve corresponder à
        entrada no arquivo \texttt{.bib} do projeto.
\end{itemize}

\section*{Seções Obrigatórias}
Cada fichamento deve conter obrigatoriamente as seguintes seções:

\begin{enumerate}
  \item \textbf{Identificação} — Dados bibliográficos completos conforme
        NBR 6023.
  \item \textbf{Objetivo} — Síntese do propósito do artigo em 3--5 linhas.
  \item \textbf{Metodologia} — Desenho do estudo, amostra, métricas e
        validação.
  \item \textbf{Resultados Principais} — Valores numéricos relevantes em
        lista itemizada.
  \item \textbf{Limitações} — Pelo menos duas limitações identificadas.
  \item \textbf{Contribuição para o Livro} — Relação explícita entre o
        artigo e o(s) capítulo(s) que ele fundamenta.
  \item \textbf{Análise Crítica} — Validade interna/externa,
        reprodutibilidade, relevância clínica, conflito de interesses e
        viés.
  \item \textbf{Citações Relevantes} — Transcrição de trechos citados no
        livro com indicação de página.
  \item \textbf{Mapeamento SDD} — Relação problema-entrada-saída-critérios
        conforme metodologia SDD do livro.
  \item \textbf{Referência ABNT} — Referência formatada conforme NBR 6023
        dentro de ambiente \texttt{verbatim}.
\end{enumerate}

\section*{Critérios de Qualidade}
\begin{itemize}
  \item \textbf{Rastreabilidade:} Toda afirmação no livro apoiada por
        \texttt{\textbackslash citacaoativa} deve ter fichamento
        correspondente.
  \item \textbf{Precisão:} Métricas numéricas (acurácia, sensibilidade,
        AUC, etc.) devem ser transcritas exatamente como publicadas.
  \item \textbf{Imparcialidade:} A análise crítica deve apontar tanto
        contribuições quanto limitações, sem viés de confirmação.
  \item \textbf{SDD:} O mapeamento SDD deve explicitar o problema
        odontológico, os tipos de dado de entrada/saída e os critérios de
        aceitação reportados no artigo.
  \item \textbf{Confiança:} Atribuir valor entre 0 (mínima) e 1 (máxima)
        refletindo a confiabilidade do estudo com base no desenho, tamanho
        amostral, validação e reprodutibilidade.
\end{itemize}

\section*{Verificação de DOIs}
Antes de finalizar o fichamento, verifique se o DOI do artigo está ativo
em \url{https://doi.org/}. DOIs que retornam \texttt{HTTP 404} ou
\texttt{HTTP 403} devem ser sinalizados no campo \textbf{DOI} com a
observação \texttt{(pendente -- possível DOI fabricado)} e reportados no
arquivo \texttt{00-doi-audit.tex}.

\section*{Integração ao Documento Principal}
Para incluir os fichamentos no livro, adicione no preâmbulo:
\begin{verbatim}
% Inclusão dos fichamentos críticos (apêndice documental)
%==============================================================================
% INSTRUÇÕES DE USO — Fichamentos Críticos
% Livro: Odontologia & Inteligência Artificial
% Conforme NBR 6023 (Referências) e NBR 10520 (Citações)
%==============================================================================

\chapter*{Instruções para Elaboração dos Fichamentos Críticos}
\addcontentsline{toc}{chapter}{Instruções — Fichamentos Críticos}

\section*{Propósito}
Os fichamentos críticos constituem o registro sistematizado da fundamentação
científica do livro. Cada artigo citado via \texttt{\textbackslash citacaoativa}
no texto deve possuir um fichamento correspondente, assegurando
rastreabilidade, auditabilidade e reprodutibilidade das fontes.

\section*{Estrutura do Arquivo}
Cada fichamento deve ser salvo como \texttt{fichamento-ODTXXX.tex}, onde
\texttt{XXX} é o identificador sequencial (001, 002, \ldots, 062).
\begin{itemize}
  \item Utilize o arquivo \texttt{template-fichamento.tex} como base.
  \item Substitua todos os campos entre colchetes pelos valores
        correspondentes do artigo fichado.
  \item A chave bibliográfica (\texttt{CHAVE\_BIB}) deve corresponder à
        entrada no arquivo \texttt{.bib} do projeto.
\end{itemize}

\section*{Seções Obrigatórias}
Cada fichamento deve conter obrigatoriamente as seguintes seções:

\begin{enumerate}
  \item \textbf{Identificação} — Dados bibliográficos completos conforme
        NBR 6023.
  \item \textbf{Objetivo} — Síntese do propósito do artigo em 3--5 linhas.
  \item \textbf{Metodologia} — Desenho do estudo, amostra, métricas e
        validação.
  \item \textbf{Resultados Principais} — Valores numéricos relevantes em
        lista itemizada.
  \item \textbf{Limitações} — Pelo menos duas limitações identificadas.
  \item \textbf{Contribuição para o Livro} — Relação explícita entre o
        artigo e o(s) capítulo(s) que ele fundamenta.
  \item \textbf{Análise Crítica} — Validade interna/externa,
        reprodutibilidade, relevância clínica, conflito de interesses e
        viés.
  \item \textbf{Citações Relevantes} — Transcrição de trechos citados no
        livro com indicação de página.
  \item \textbf{Mapeamento SDD} — Relação problema-entrada-saída-critérios
        conforme metodologia SDD do livro.
  \item \textbf{Referência ABNT} — Referência formatada conforme NBR 6023
        dentro de ambiente \texttt{verbatim}.
\end{enumerate}

\section*{Critérios de Qualidade}
\begin{itemize}
  \item \textbf{Rastreabilidade:} Toda afirmação no livro apoiada por
        \texttt{\textbackslash citacaoativa} deve ter fichamento
        correspondente.
  \item \textbf{Precisão:} Métricas numéricas (acurácia, sensibilidade,
        AUC, etc.) devem ser transcritas exatamente como publicadas.
  \item \textbf{Imparcialidade:} A análise crítica deve apontar tanto
        contribuições quanto limitações, sem viés de confirmação.
  \item \textbf{SDD:} O mapeamento SDD deve explicitar o problema
        odontológico, os tipos de dado de entrada/saída e os critérios de
        aceitação reportados no artigo.
  \item \textbf{Confiança:} Atribuir valor entre 0 (mínima) e 1 (máxima)
        refletindo a confiabilidade do estudo com base no desenho, tamanho
        amostral, validação e reprodutibilidade.
\end{itemize}

\section*{Verificação de DOIs}
Antes de finalizar o fichamento, verifique se o DOI do artigo está ativo
em \url{https://doi.org/}. DOIs que retornam \texttt{HTTP 404} ou
\texttt{HTTP 403} devem ser sinalizados no campo \textbf{DOI} com a
observação \texttt{(pendente -- possível DOI fabricado)} e reportados no
arquivo \texttt{00-doi-audit.tex}.

\section*{Integração ao Documento Principal}
Para incluir os fichamentos no livro, adicione no preâmbulo:
\begin{verbatim}
% Inclusão dos fichamentos críticos (apêndice documental)
%==============================================================================
% INSTRUÇÕES DE USO — Fichamentos Críticos
% Livro: Odontologia & Inteligência Artificial
% Conforme NBR 6023 (Referências) e NBR 10520 (Citações)
%==============================================================================

\chapter*{Instruções para Elaboração dos Fichamentos Críticos}
\addcontentsline{toc}{chapter}{Instruções — Fichamentos Críticos}

\section*{Propósito}
Os fichamentos críticos constituem o registro sistematizado da fundamentação
científica do livro. Cada artigo citado via \texttt{\textbackslash citacaoativa}
no texto deve possuir um fichamento correspondente, assegurando
rastreabilidade, auditabilidade e reprodutibilidade das fontes.

\section*{Estrutura do Arquivo}
Cada fichamento deve ser salvo como \texttt{fichamento-ODTXXX.tex}, onde
\texttt{XXX} é o identificador sequencial (001, 002, \ldots, 062).
\begin{itemize}
  \item Utilize o arquivo \texttt{template-fichamento.tex} como base.
  \item Substitua todos os campos entre colchetes pelos valores
        correspondentes do artigo fichado.
  \item A chave bibliográfica (\texttt{CHAVE\_BIB}) deve corresponder à
        entrada no arquivo \texttt{.bib} do projeto.
\end{itemize}

\section*{Seções Obrigatórias}
Cada fichamento deve conter obrigatoriamente as seguintes seções:

\begin{enumerate}
  \item \textbf{Identificação} — Dados bibliográficos completos conforme
        NBR 6023.
  \item \textbf{Objetivo} — Síntese do propósito do artigo em 3--5 linhas.
  \item \textbf{Metodologia} — Desenho do estudo, amostra, métricas e
        validação.
  \item \textbf{Resultados Principais} — Valores numéricos relevantes em
        lista itemizada.
  \item \textbf{Limitações} — Pelo menos duas limitações identificadas.
  \item \textbf{Contribuição para o Livro} — Relação explícita entre o
        artigo e o(s) capítulo(s) que ele fundamenta.
  \item \textbf{Análise Crítica} — Validade interna/externa,
        reprodutibilidade, relevância clínica, conflito de interesses e
        viés.
  \item \textbf{Citações Relevantes} — Transcrição de trechos citados no
        livro com indicação de página.
  \item \textbf{Mapeamento SDD} — Relação problema-entrada-saída-critérios
        conforme metodologia SDD do livro.
  \item \textbf{Referência ABNT} — Referência formatada conforme NBR 6023
        dentro de ambiente \texttt{verbatim}.
\end{enumerate}

\section*{Critérios de Qualidade}
\begin{itemize}
  \item \textbf{Rastreabilidade:} Toda afirmação no livro apoiada por
        \texttt{\textbackslash citacaoativa} deve ter fichamento
        correspondente.
  \item \textbf{Precisão:} Métricas numéricas (acurácia, sensibilidade,
        AUC, etc.) devem ser transcritas exatamente como publicadas.
  \item \textbf{Imparcialidade:} A análise crítica deve apontar tanto
        contribuições quanto limitações, sem viés de confirmação.
  \item \textbf{SDD:} O mapeamento SDD deve explicitar o problema
        odontológico, os tipos de dado de entrada/saída e os critérios de
        aceitação reportados no artigo.
  \item \textbf{Confiança:} Atribuir valor entre 0 (mínima) e 1 (máxima)
        refletindo a confiabilidade do estudo com base no desenho, tamanho
        amostral, validação e reprodutibilidade.
\end{itemize}

\section*{Verificação de DOIs}
Antes de finalizar o fichamento, verifique se o DOI do artigo está ativo
em \url{https://doi.org/}. DOIs que retornam \texttt{HTTP 404} ou
\texttt{HTTP 403} devem ser sinalizados no campo \textbf{DOI} com a
observação \texttt{(pendente -- possível DOI fabricado)} e reportados no
arquivo \texttt{00-doi-audit.tex}.

\section*{Integração ao Documento Principal}
Para incluir os fichamentos no livro, adicione no preâmbulo:
\begin{verbatim}
% Inclusão dos fichamentos críticos (apêndice documental)
\input{fichamentos/00-instructions}
\input{fichamentos/fichamento-ODT001}
\input{fichamentos/fichamento-ODT002}
% ...
\end{verbatim}
Ou, alternativamente, crie um apêndice específico:
\begin{verbatim}
\begin{apendicesenv}
\partapendices
\chapter{FICHAMENTOS CRÍTICOS}
\input{fichamentos/00-instructions}
\input{fichamentos/fichamento-ODT001}
% ...
\end{apendicesenv}
\end{verbatim}

\section*{Exemplo de Arquivo}
Consulte \texttt{template-fichamento.tex} para o modelo completo e
\texttt{00-doi-audit.tex} para o relatório de auditoria de DOIs.

\vfill
\begin{center}
\rule{0.5\textwidth}{0.4pt}\\
\small{Documento elaborado conforme NBR 14724 e NBR 6029.}\\
\small{Revisão Qualis A1.}
\end{center}

\input{fichamentos/fichamento-ODT001}
\input{fichamentos/fichamento-ODT002}
% ...
\end{verbatim}
Ou, alternativamente, crie um apêndice específico:
\begin{verbatim}
\begin{apendicesenv}
\partapendices
\chapter{FICHAMENTOS CRÍTICOS}
%==============================================================================
% INSTRUÇÕES DE USO — Fichamentos Críticos
% Livro: Odontologia & Inteligência Artificial
% Conforme NBR 6023 (Referências) e NBR 10520 (Citações)
%==============================================================================

\chapter*{Instruções para Elaboração dos Fichamentos Críticos}
\addcontentsline{toc}{chapter}{Instruções — Fichamentos Críticos}

\section*{Propósito}
Os fichamentos críticos constituem o registro sistematizado da fundamentação
científica do livro. Cada artigo citado via \texttt{\textbackslash citacaoativa}
no texto deve possuir um fichamento correspondente, assegurando
rastreabilidade, auditabilidade e reprodutibilidade das fontes.

\section*{Estrutura do Arquivo}
Cada fichamento deve ser salvo como \texttt{fichamento-ODTXXX.tex}, onde
\texttt{XXX} é o identificador sequencial (001, 002, \ldots, 062).
\begin{itemize}
  \item Utilize o arquivo \texttt{template-fichamento.tex} como base.
  \item Substitua todos os campos entre colchetes pelos valores
        correspondentes do artigo fichado.
  \item A chave bibliográfica (\texttt{CHAVE\_BIB}) deve corresponder à
        entrada no arquivo \texttt{.bib} do projeto.
\end{itemize}

\section*{Seções Obrigatórias}
Cada fichamento deve conter obrigatoriamente as seguintes seções:

\begin{enumerate}
  \item \textbf{Identificação} — Dados bibliográficos completos conforme
        NBR 6023.
  \item \textbf{Objetivo} — Síntese do propósito do artigo em 3--5 linhas.
  \item \textbf{Metodologia} — Desenho do estudo, amostra, métricas e
        validação.
  \item \textbf{Resultados Principais} — Valores numéricos relevantes em
        lista itemizada.
  \item \textbf{Limitações} — Pelo menos duas limitações identificadas.
  \item \textbf{Contribuição para o Livro} — Relação explícita entre o
        artigo e o(s) capítulo(s) que ele fundamenta.
  \item \textbf{Análise Crítica} — Validade interna/externa,
        reprodutibilidade, relevância clínica, conflito de interesses e
        viés.
  \item \textbf{Citações Relevantes} — Transcrição de trechos citados no
        livro com indicação de página.
  \item \textbf{Mapeamento SDD} — Relação problema-entrada-saída-critérios
        conforme metodologia SDD do livro.
  \item \textbf{Referência ABNT} — Referência formatada conforme NBR 6023
        dentro de ambiente \texttt{verbatim}.
\end{enumerate}

\section*{Critérios de Qualidade}
\begin{itemize}
  \item \textbf{Rastreabilidade:} Toda afirmação no livro apoiada por
        \texttt{\textbackslash citacaoativa} deve ter fichamento
        correspondente.
  \item \textbf{Precisão:} Métricas numéricas (acurácia, sensibilidade,
        AUC, etc.) devem ser transcritas exatamente como publicadas.
  \item \textbf{Imparcialidade:} A análise crítica deve apontar tanto
        contribuições quanto limitações, sem viés de confirmação.
  \item \textbf{SDD:} O mapeamento SDD deve explicitar o problema
        odontológico, os tipos de dado de entrada/saída e os critérios de
        aceitação reportados no artigo.
  \item \textbf{Confiança:} Atribuir valor entre 0 (mínima) e 1 (máxima)
        refletindo a confiabilidade do estudo com base no desenho, tamanho
        amostral, validação e reprodutibilidade.
\end{itemize}

\section*{Verificação de DOIs}
Antes de finalizar o fichamento, verifique se o DOI do artigo está ativo
em \url{https://doi.org/}. DOIs que retornam \texttt{HTTP 404} ou
\texttt{HTTP 403} devem ser sinalizados no campo \textbf{DOI} com a
observação \texttt{(pendente -- possível DOI fabricado)} e reportados no
arquivo \texttt{00-doi-audit.tex}.

\section*{Integração ao Documento Principal}
Para incluir os fichamentos no livro, adicione no preâmbulo:
\begin{verbatim}
% Inclusão dos fichamentos críticos (apêndice documental)
\input{fichamentos/00-instructions}
\input{fichamentos/fichamento-ODT001}
\input{fichamentos/fichamento-ODT002}
% ...
\end{verbatim}
Ou, alternativamente, crie um apêndice específico:
\begin{verbatim}
\begin{apendicesenv}
\partapendices
\chapter{FICHAMENTOS CRÍTICOS}
\input{fichamentos/00-instructions}
\input{fichamentos/fichamento-ODT001}
% ...
\end{apendicesenv}
\end{verbatim}

\section*{Exemplo de Arquivo}
Consulte \texttt{template-fichamento.tex} para o modelo completo e
\texttt{00-doi-audit.tex} para o relatório de auditoria de DOIs.

\vfill
\begin{center}
\rule{0.5\textwidth}{0.4pt}\\
\small{Documento elaborado conforme NBR 14724 e NBR 6029.}\\
\small{Revisão Qualis A1.}
\end{center}

\input{fichamentos/fichamento-ODT001}
% ...
\end{apendicesenv}
\end{verbatim}

\section*{Exemplo de Arquivo}
Consulte \texttt{template-fichamento.tex} para o modelo completo e
\texttt{00-doi-audit.tex} para o relatório de auditoria de DOIs.

\vfill
\begin{center}
\rule{0.5\textwidth}{0.4pt}\\
\small{Documento elaborado conforme NBR 14724 e NBR 6029.}\\
\small{Revisão Qualis A1.}
\end{center}

\input{fichamentos/fichamento-ODT001}
\input{fichamentos/fichamento-ODT002}
% ...
\end{verbatim}
Ou, alternativamente, crie um apêndice específico:
\begin{verbatim}
\begin{apendicesenv}
\partapendices
\chapter{FICHAMENTOS CRÍTICOS}
%==============================================================================
% INSTRUÇÕES DE USO — Fichamentos Críticos
% Livro: Odontologia & Inteligência Artificial
% Conforme NBR 6023 (Referências) e NBR 10520 (Citações)
%==============================================================================

\chapter*{Instruções para Elaboração dos Fichamentos Críticos}
\addcontentsline{toc}{chapter}{Instruções — Fichamentos Críticos}

\section*{Propósito}
Os fichamentos críticos constituem o registro sistematizado da fundamentação
científica do livro. Cada artigo citado via \texttt{\textbackslash citacaoativa}
no texto deve possuir um fichamento correspondente, assegurando
rastreabilidade, auditabilidade e reprodutibilidade das fontes.

\section*{Estrutura do Arquivo}
Cada fichamento deve ser salvo como \texttt{fichamento-ODTXXX.tex}, onde
\texttt{XXX} é o identificador sequencial (001, 002, \ldots, 062).
\begin{itemize}
  \item Utilize o arquivo \texttt{template-fichamento.tex} como base.
  \item Substitua todos os campos entre colchetes pelos valores
        correspondentes do artigo fichado.
  \item A chave bibliográfica (\texttt{CHAVE\_BIB}) deve corresponder à
        entrada no arquivo \texttt{.bib} do projeto.
\end{itemize}

\section*{Seções Obrigatórias}
Cada fichamento deve conter obrigatoriamente as seguintes seções:

\begin{enumerate}
  \item \textbf{Identificação} — Dados bibliográficos completos conforme
        NBR 6023.
  \item \textbf{Objetivo} — Síntese do propósito do artigo em 3--5 linhas.
  \item \textbf{Metodologia} — Desenho do estudo, amostra, métricas e
        validação.
  \item \textbf{Resultados Principais} — Valores numéricos relevantes em
        lista itemizada.
  \item \textbf{Limitações} — Pelo menos duas limitações identificadas.
  \item \textbf{Contribuição para o Livro} — Relação explícita entre o
        artigo e o(s) capítulo(s) que ele fundamenta.
  \item \textbf{Análise Crítica} — Validade interna/externa,
        reprodutibilidade, relevância clínica, conflito de interesses e
        viés.
  \item \textbf{Citações Relevantes} — Transcrição de trechos citados no
        livro com indicação de página.
  \item \textbf{Mapeamento SDD} — Relação problema-entrada-saída-critérios
        conforme metodologia SDD do livro.
  \item \textbf{Referência ABNT} — Referência formatada conforme NBR 6023
        dentro de ambiente \texttt{verbatim}.
\end{enumerate}

\section*{Critérios de Qualidade}
\begin{itemize}
  \item \textbf{Rastreabilidade:} Toda afirmação no livro apoiada por
        \texttt{\textbackslash citacaoativa} deve ter fichamento
        correspondente.
  \item \textbf{Precisão:} Métricas numéricas (acurácia, sensibilidade,
        AUC, etc.) devem ser transcritas exatamente como publicadas.
  \item \textbf{Imparcialidade:} A análise crítica deve apontar tanto
        contribuições quanto limitações, sem viés de confirmação.
  \item \textbf{SDD:} O mapeamento SDD deve explicitar o problema
        odontológico, os tipos de dado de entrada/saída e os critérios de
        aceitação reportados no artigo.
  \item \textbf{Confiança:} Atribuir valor entre 0 (mínima) e 1 (máxima)
        refletindo a confiabilidade do estudo com base no desenho, tamanho
        amostral, validação e reprodutibilidade.
\end{itemize}

\section*{Verificação de DOIs}
Antes de finalizar o fichamento, verifique se o DOI do artigo está ativo
em \url{https://doi.org/}. DOIs que retornam \texttt{HTTP 404} ou
\texttt{HTTP 403} devem ser sinalizados no campo \textbf{DOI} com a
observação \texttt{(pendente -- possível DOI fabricado)} e reportados no
arquivo \texttt{00-doi-audit.tex}.

\section*{Integração ao Documento Principal}
Para incluir os fichamentos no livro, adicione no preâmbulo:
\begin{verbatim}
% Inclusão dos fichamentos críticos (apêndice documental)
%==============================================================================
% INSTRUÇÕES DE USO — Fichamentos Críticos
% Livro: Odontologia & Inteligência Artificial
% Conforme NBR 6023 (Referências) e NBR 10520 (Citações)
%==============================================================================

\chapter*{Instruções para Elaboração dos Fichamentos Críticos}
\addcontentsline{toc}{chapter}{Instruções — Fichamentos Críticos}

\section*{Propósito}
Os fichamentos críticos constituem o registro sistematizado da fundamentação
científica do livro. Cada artigo citado via \texttt{\textbackslash citacaoativa}
no texto deve possuir um fichamento correspondente, assegurando
rastreabilidade, auditabilidade e reprodutibilidade das fontes.

\section*{Estrutura do Arquivo}
Cada fichamento deve ser salvo como \texttt{fichamento-ODTXXX.tex}, onde
\texttt{XXX} é o identificador sequencial (001, 002, \ldots, 062).
\begin{itemize}
  \item Utilize o arquivo \texttt{template-fichamento.tex} como base.
  \item Substitua todos os campos entre colchetes pelos valores
        correspondentes do artigo fichado.
  \item A chave bibliográfica (\texttt{CHAVE\_BIB}) deve corresponder à
        entrada no arquivo \texttt{.bib} do projeto.
\end{itemize}

\section*{Seções Obrigatórias}
Cada fichamento deve conter obrigatoriamente as seguintes seções:

\begin{enumerate}
  \item \textbf{Identificação} — Dados bibliográficos completos conforme
        NBR 6023.
  \item \textbf{Objetivo} — Síntese do propósito do artigo em 3--5 linhas.
  \item \textbf{Metodologia} — Desenho do estudo, amostra, métricas e
        validação.
  \item \textbf{Resultados Principais} — Valores numéricos relevantes em
        lista itemizada.
  \item \textbf{Limitações} — Pelo menos duas limitações identificadas.
  \item \textbf{Contribuição para o Livro} — Relação explícita entre o
        artigo e o(s) capítulo(s) que ele fundamenta.
  \item \textbf{Análise Crítica} — Validade interna/externa,
        reprodutibilidade, relevância clínica, conflito de interesses e
        viés.
  \item \textbf{Citações Relevantes} — Transcrição de trechos citados no
        livro com indicação de página.
  \item \textbf{Mapeamento SDD} — Relação problema-entrada-saída-critérios
        conforme metodologia SDD do livro.
  \item \textbf{Referência ABNT} — Referência formatada conforme NBR 6023
        dentro de ambiente \texttt{verbatim}.
\end{enumerate}

\section*{Critérios de Qualidade}
\begin{itemize}
  \item \textbf{Rastreabilidade:} Toda afirmação no livro apoiada por
        \texttt{\textbackslash citacaoativa} deve ter fichamento
        correspondente.
  \item \textbf{Precisão:} Métricas numéricas (acurácia, sensibilidade,
        AUC, etc.) devem ser transcritas exatamente como publicadas.
  \item \textbf{Imparcialidade:} A análise crítica deve apontar tanto
        contribuições quanto limitações, sem viés de confirmação.
  \item \textbf{SDD:} O mapeamento SDD deve explicitar o problema
        odontológico, os tipos de dado de entrada/saída e os critérios de
        aceitação reportados no artigo.
  \item \textbf{Confiança:} Atribuir valor entre 0 (mínima) e 1 (máxima)
        refletindo a confiabilidade do estudo com base no desenho, tamanho
        amostral, validação e reprodutibilidade.
\end{itemize}

\section*{Verificação de DOIs}
Antes de finalizar o fichamento, verifique se o DOI do artigo está ativo
em \url{https://doi.org/}. DOIs que retornam \texttt{HTTP 404} ou
\texttt{HTTP 403} devem ser sinalizados no campo \textbf{DOI} com a
observação \texttt{(pendente -- possível DOI fabricado)} e reportados no
arquivo \texttt{00-doi-audit.tex}.

\section*{Integração ao Documento Principal}
Para incluir os fichamentos no livro, adicione no preâmbulo:
\begin{verbatim}
% Inclusão dos fichamentos críticos (apêndice documental)
\input{fichamentos/00-instructions}
\input{fichamentos/fichamento-ODT001}
\input{fichamentos/fichamento-ODT002}
% ...
\end{verbatim}
Ou, alternativamente, crie um apêndice específico:
\begin{verbatim}
\begin{apendicesenv}
\partapendices
\chapter{FICHAMENTOS CRÍTICOS}
\input{fichamentos/00-instructions}
\input{fichamentos/fichamento-ODT001}
% ...
\end{apendicesenv}
\end{verbatim}

\section*{Exemplo de Arquivo}
Consulte \texttt{template-fichamento.tex} para o modelo completo e
\texttt{00-doi-audit.tex} para o relatório de auditoria de DOIs.

\vfill
\begin{center}
\rule{0.5\textwidth}{0.4pt}\\
\small{Documento elaborado conforme NBR 14724 e NBR 6029.}\\
\small{Revisão Qualis A1.}
\end{center}

\input{fichamentos/fichamento-ODT001}
\input{fichamentos/fichamento-ODT002}
% ...
\end{verbatim}
Ou, alternativamente, crie um apêndice específico:
\begin{verbatim}
\begin{apendicesenv}
\partapendices
\chapter{FICHAMENTOS CRÍTICOS}
%==============================================================================
% INSTRUÇÕES DE USO — Fichamentos Críticos
% Livro: Odontologia & Inteligência Artificial
% Conforme NBR 6023 (Referências) e NBR 10520 (Citações)
%==============================================================================

\chapter*{Instruções para Elaboração dos Fichamentos Críticos}
\addcontentsline{toc}{chapter}{Instruções — Fichamentos Críticos}

\section*{Propósito}
Os fichamentos críticos constituem o registro sistematizado da fundamentação
científica do livro. Cada artigo citado via \texttt{\textbackslash citacaoativa}
no texto deve possuir um fichamento correspondente, assegurando
rastreabilidade, auditabilidade e reprodutibilidade das fontes.

\section*{Estrutura do Arquivo}
Cada fichamento deve ser salvo como \texttt{fichamento-ODTXXX.tex}, onde
\texttt{XXX} é o identificador sequencial (001, 002, \ldots, 062).
\begin{itemize}
  \item Utilize o arquivo \texttt{template-fichamento.tex} como base.
  \item Substitua todos os campos entre colchetes pelos valores
        correspondentes do artigo fichado.
  \item A chave bibliográfica (\texttt{CHAVE\_BIB}) deve corresponder à
        entrada no arquivo \texttt{.bib} do projeto.
\end{itemize}

\section*{Seções Obrigatórias}
Cada fichamento deve conter obrigatoriamente as seguintes seções:

\begin{enumerate}
  \item \textbf{Identificação} — Dados bibliográficos completos conforme
        NBR 6023.
  \item \textbf{Objetivo} — Síntese do propósito do artigo em 3--5 linhas.
  \item \textbf{Metodologia} — Desenho do estudo, amostra, métricas e
        validação.
  \item \textbf{Resultados Principais} — Valores numéricos relevantes em
        lista itemizada.
  \item \textbf{Limitações} — Pelo menos duas limitações identificadas.
  \item \textbf{Contribuição para o Livro} — Relação explícita entre o
        artigo e o(s) capítulo(s) que ele fundamenta.
  \item \textbf{Análise Crítica} — Validade interna/externa,
        reprodutibilidade, relevância clínica, conflito de interesses e
        viés.
  \item \textbf{Citações Relevantes} — Transcrição de trechos citados no
        livro com indicação de página.
  \item \textbf{Mapeamento SDD} — Relação problema-entrada-saída-critérios
        conforme metodologia SDD do livro.
  \item \textbf{Referência ABNT} — Referência formatada conforme NBR 6023
        dentro de ambiente \texttt{verbatim}.
\end{enumerate}

\section*{Critérios de Qualidade}
\begin{itemize}
  \item \textbf{Rastreabilidade:} Toda afirmação no livro apoiada por
        \texttt{\textbackslash citacaoativa} deve ter fichamento
        correspondente.
  \item \textbf{Precisão:} Métricas numéricas (acurácia, sensibilidade,
        AUC, etc.) devem ser transcritas exatamente como publicadas.
  \item \textbf{Imparcialidade:} A análise crítica deve apontar tanto
        contribuições quanto limitações, sem viés de confirmação.
  \item \textbf{SDD:} O mapeamento SDD deve explicitar o problema
        odontológico, os tipos de dado de entrada/saída e os critérios de
        aceitação reportados no artigo.
  \item \textbf{Confiança:} Atribuir valor entre 0 (mínima) e 1 (máxima)
        refletindo a confiabilidade do estudo com base no desenho, tamanho
        amostral, validação e reprodutibilidade.
\end{itemize}

\section*{Verificação de DOIs}
Antes de finalizar o fichamento, verifique se o DOI do artigo está ativo
em \url{https://doi.org/}. DOIs que retornam \texttt{HTTP 404} ou
\texttt{HTTP 403} devem ser sinalizados no campo \textbf{DOI} com a
observação \texttt{(pendente -- possível DOI fabricado)} e reportados no
arquivo \texttt{00-doi-audit.tex}.

\section*{Integração ao Documento Principal}
Para incluir os fichamentos no livro, adicione no preâmbulo:
\begin{verbatim}
% Inclusão dos fichamentos críticos (apêndice documental)
\input{fichamentos/00-instructions}
\input{fichamentos/fichamento-ODT001}
\input{fichamentos/fichamento-ODT002}
% ...
\end{verbatim}
Ou, alternativamente, crie um apêndice específico:
\begin{verbatim}
\begin{apendicesenv}
\partapendices
\chapter{FICHAMENTOS CRÍTICOS}
\input{fichamentos/00-instructions}
\input{fichamentos/fichamento-ODT001}
% ...
\end{apendicesenv}
\end{verbatim}

\section*{Exemplo de Arquivo}
Consulte \texttt{template-fichamento.tex} para o modelo completo e
\texttt{00-doi-audit.tex} para o relatório de auditoria de DOIs.

\vfill
\begin{center}
\rule{0.5\textwidth}{0.4pt}\\
\small{Documento elaborado conforme NBR 14724 e NBR 6029.}\\
\small{Revisão Qualis A1.}
\end{center}

\input{fichamentos/fichamento-ODT001}
% ...
\end{apendicesenv}
\end{verbatim}

\section*{Exemplo de Arquivo}
Consulte \texttt{template-fichamento.tex} para o modelo completo e
\texttt{00-doi-audit.tex} para o relatório de auditoria de DOIs.

\vfill
\begin{center}
\rule{0.5\textwidth}{0.4pt}\\
\small{Documento elaborado conforme NBR 14724 e NBR 6029.}\\
\small{Revisão Qualis A1.}
\end{center}

\input{fichamentos/fichamento-ODT001}
% ...
\end{apendicesenv}
\end{verbatim}

\section*{Exemplo de Arquivo}
Consulte \texttt{template-fichamento.tex} para o modelo completo e
\texttt{00-doi-audit.tex} para o relatório de auditoria de DOIs.

\vfill
\begin{center}
\rule{0.5\textwidth}{0.4pt}\\
\small{Documento elaborado conforme NBR 14724 e NBR 6029.}\\
\small{Revisão Qualis A1.}
\end{center}

\input{fichamentos/fichamento-ODT001}
\input{fichamentos/fichamento-ODT002}
% ...
\end{verbatim}
Ou, alternativamente, crie um apêndice específico:
\begin{verbatim}
\begin{apendicesenv}
\partapendices
\chapter{FICHAMENTOS CRÍTICOS}
%==============================================================================
% INSTRUÇÕES DE USO — Fichamentos Críticos
% Livro: Odontologia & Inteligência Artificial
% Conforme NBR 6023 (Referências) e NBR 10520 (Citações)
%==============================================================================

\chapter*{Instruções para Elaboração dos Fichamentos Críticos}
\addcontentsline{toc}{chapter}{Instruções — Fichamentos Críticos}

\section*{Propósito}
Os fichamentos críticos constituem o registro sistematizado da fundamentação
científica do livro. Cada artigo citado via \texttt{\textbackslash citacaoativa}
no texto deve possuir um fichamento correspondente, assegurando
rastreabilidade, auditabilidade e reprodutibilidade das fontes.

\section*{Estrutura do Arquivo}
Cada fichamento deve ser salvo como \texttt{fichamento-ODTXXX.tex}, onde
\texttt{XXX} é o identificador sequencial (001, 002, \ldots, 062).
\begin{itemize}
  \item Utilize o arquivo \texttt{template-fichamento.tex} como base.
  \item Substitua todos os campos entre colchetes pelos valores
        correspondentes do artigo fichado.
  \item A chave bibliográfica (\texttt{CHAVE\_BIB}) deve corresponder à
        entrada no arquivo \texttt{.bib} do projeto.
\end{itemize}

\section*{Seções Obrigatórias}
Cada fichamento deve conter obrigatoriamente as seguintes seções:

\begin{enumerate}
  \item \textbf{Identificação} — Dados bibliográficos completos conforme
        NBR 6023.
  \item \textbf{Objetivo} — Síntese do propósito do artigo em 3--5 linhas.
  \item \textbf{Metodologia} — Desenho do estudo, amostra, métricas e
        validação.
  \item \textbf{Resultados Principais} — Valores numéricos relevantes em
        lista itemizada.
  \item \textbf{Limitações} — Pelo menos duas limitações identificadas.
  \item \textbf{Contribuição para o Livro} — Relação explícita entre o
        artigo e o(s) capítulo(s) que ele fundamenta.
  \item \textbf{Análise Crítica} — Validade interna/externa,
        reprodutibilidade, relevância clínica, conflito de interesses e
        viés.
  \item \textbf{Citações Relevantes} — Transcrição de trechos citados no
        livro com indicação de página.
  \item \textbf{Mapeamento SDD} — Relação problema-entrada-saída-critérios
        conforme metodologia SDD do livro.
  \item \textbf{Referência ABNT} — Referência formatada conforme NBR 6023
        dentro de ambiente \texttt{verbatim}.
\end{enumerate}

\section*{Critérios de Qualidade}
\begin{itemize}
  \item \textbf{Rastreabilidade:} Toda afirmação no livro apoiada por
        \texttt{\textbackslash citacaoativa} deve ter fichamento
        correspondente.
  \item \textbf{Precisão:} Métricas numéricas (acurácia, sensibilidade,
        AUC, etc.) devem ser transcritas exatamente como publicadas.
  \item \textbf{Imparcialidade:} A análise crítica deve apontar tanto
        contribuições quanto limitações, sem viés de confirmação.
  \item \textbf{SDD:} O mapeamento SDD deve explicitar o problema
        odontológico, os tipos de dado de entrada/saída e os critérios de
        aceitação reportados no artigo.
  \item \textbf{Confiança:} Atribuir valor entre 0 (mínima) e 1 (máxima)
        refletindo a confiabilidade do estudo com base no desenho, tamanho
        amostral, validação e reprodutibilidade.
\end{itemize}

\section*{Verificação de DOIs}
Antes de finalizar o fichamento, verifique se o DOI do artigo está ativo
em \url{https://doi.org/}. DOIs que retornam \texttt{HTTP 404} ou
\texttt{HTTP 403} devem ser sinalizados no campo \textbf{DOI} com a
observação \texttt{(pendente -- possível DOI fabricado)} e reportados no
arquivo \texttt{00-doi-audit.tex}.

\section*{Integração ao Documento Principal}
Para incluir os fichamentos no livro, adicione no preâmbulo:
\begin{verbatim}
% Inclusão dos fichamentos críticos (apêndice documental)
%==============================================================================
% INSTRUÇÕES DE USO — Fichamentos Críticos
% Livro: Odontologia & Inteligência Artificial
% Conforme NBR 6023 (Referências) e NBR 10520 (Citações)
%==============================================================================

\chapter*{Instruções para Elaboração dos Fichamentos Críticos}
\addcontentsline{toc}{chapter}{Instruções — Fichamentos Críticos}

\section*{Propósito}
Os fichamentos críticos constituem o registro sistematizado da fundamentação
científica do livro. Cada artigo citado via \texttt{\textbackslash citacaoativa}
no texto deve possuir um fichamento correspondente, assegurando
rastreabilidade, auditabilidade e reprodutibilidade das fontes.

\section*{Estrutura do Arquivo}
Cada fichamento deve ser salvo como \texttt{fichamento-ODTXXX.tex}, onde
\texttt{XXX} é o identificador sequencial (001, 002, \ldots, 062).
\begin{itemize}
  \item Utilize o arquivo \texttt{template-fichamento.tex} como base.
  \item Substitua todos os campos entre colchetes pelos valores
        correspondentes do artigo fichado.
  \item A chave bibliográfica (\texttt{CHAVE\_BIB}) deve corresponder à
        entrada no arquivo \texttt{.bib} do projeto.
\end{itemize}

\section*{Seções Obrigatórias}
Cada fichamento deve conter obrigatoriamente as seguintes seções:

\begin{enumerate}
  \item \textbf{Identificação} — Dados bibliográficos completos conforme
        NBR 6023.
  \item \textbf{Objetivo} — Síntese do propósito do artigo em 3--5 linhas.
  \item \textbf{Metodologia} — Desenho do estudo, amostra, métricas e
        validação.
  \item \textbf{Resultados Principais} — Valores numéricos relevantes em
        lista itemizada.
  \item \textbf{Limitações} — Pelo menos duas limitações identificadas.
  \item \textbf{Contribuição para o Livro} — Relação explícita entre o
        artigo e o(s) capítulo(s) que ele fundamenta.
  \item \textbf{Análise Crítica} — Validade interna/externa,
        reprodutibilidade, relevância clínica, conflito de interesses e
        viés.
  \item \textbf{Citações Relevantes} — Transcrição de trechos citados no
        livro com indicação de página.
  \item \textbf{Mapeamento SDD} — Relação problema-entrada-saída-critérios
        conforme metodologia SDD do livro.
  \item \textbf{Referência ABNT} — Referência formatada conforme NBR 6023
        dentro de ambiente \texttt{verbatim}.
\end{enumerate}

\section*{Critérios de Qualidade}
\begin{itemize}
  \item \textbf{Rastreabilidade:} Toda afirmação no livro apoiada por
        \texttt{\textbackslash citacaoativa} deve ter fichamento
        correspondente.
  \item \textbf{Precisão:} Métricas numéricas (acurácia, sensibilidade,
        AUC, etc.) devem ser transcritas exatamente como publicadas.
  \item \textbf{Imparcialidade:} A análise crítica deve apontar tanto
        contribuições quanto limitações, sem viés de confirmação.
  \item \textbf{SDD:} O mapeamento SDD deve explicitar o problema
        odontológico, os tipos de dado de entrada/saída e os critérios de
        aceitação reportados no artigo.
  \item \textbf{Confiança:} Atribuir valor entre 0 (mínima) e 1 (máxima)
        refletindo a confiabilidade do estudo com base no desenho, tamanho
        amostral, validação e reprodutibilidade.
\end{itemize}

\section*{Verificação de DOIs}
Antes de finalizar o fichamento, verifique se o DOI do artigo está ativo
em \url{https://doi.org/}. DOIs que retornam \texttt{HTTP 404} ou
\texttt{HTTP 403} devem ser sinalizados no campo \textbf{DOI} com a
observação \texttt{(pendente -- possível DOI fabricado)} e reportados no
arquivo \texttt{00-doi-audit.tex}.

\section*{Integração ao Documento Principal}
Para incluir os fichamentos no livro, adicione no preâmbulo:
\begin{verbatim}
% Inclusão dos fichamentos críticos (apêndice documental)
%==============================================================================
% INSTRUÇÕES DE USO — Fichamentos Críticos
% Livro: Odontologia & Inteligência Artificial
% Conforme NBR 6023 (Referências) e NBR 10520 (Citações)
%==============================================================================

\chapter*{Instruções para Elaboração dos Fichamentos Críticos}
\addcontentsline{toc}{chapter}{Instruções — Fichamentos Críticos}

\section*{Propósito}
Os fichamentos críticos constituem o registro sistematizado da fundamentação
científica do livro. Cada artigo citado via \texttt{\textbackslash citacaoativa}
no texto deve possuir um fichamento correspondente, assegurando
rastreabilidade, auditabilidade e reprodutibilidade das fontes.

\section*{Estrutura do Arquivo}
Cada fichamento deve ser salvo como \texttt{fichamento-ODTXXX.tex}, onde
\texttt{XXX} é o identificador sequencial (001, 002, \ldots, 062).
\begin{itemize}
  \item Utilize o arquivo \texttt{template-fichamento.tex} como base.
  \item Substitua todos os campos entre colchetes pelos valores
        correspondentes do artigo fichado.
  \item A chave bibliográfica (\texttt{CHAVE\_BIB}) deve corresponder à
        entrada no arquivo \texttt{.bib} do projeto.
\end{itemize}

\section*{Seções Obrigatórias}
Cada fichamento deve conter obrigatoriamente as seguintes seções:

\begin{enumerate}
  \item \textbf{Identificação} — Dados bibliográficos completos conforme
        NBR 6023.
  \item \textbf{Objetivo} — Síntese do propósito do artigo em 3--5 linhas.
  \item \textbf{Metodologia} — Desenho do estudo, amostra, métricas e
        validação.
  \item \textbf{Resultados Principais} — Valores numéricos relevantes em
        lista itemizada.
  \item \textbf{Limitações} — Pelo menos duas limitações identificadas.
  \item \textbf{Contribuição para o Livro} — Relação explícita entre o
        artigo e o(s) capítulo(s) que ele fundamenta.
  \item \textbf{Análise Crítica} — Validade interna/externa,
        reprodutibilidade, relevância clínica, conflito de interesses e
        viés.
  \item \textbf{Citações Relevantes} — Transcrição de trechos citados no
        livro com indicação de página.
  \item \textbf{Mapeamento SDD} — Relação problema-entrada-saída-critérios
        conforme metodologia SDD do livro.
  \item \textbf{Referência ABNT} — Referência formatada conforme NBR 6023
        dentro de ambiente \texttt{verbatim}.
\end{enumerate}

\section*{Critérios de Qualidade}
\begin{itemize}
  \item \textbf{Rastreabilidade:} Toda afirmação no livro apoiada por
        \texttt{\textbackslash citacaoativa} deve ter fichamento
        correspondente.
  \item \textbf{Precisão:} Métricas numéricas (acurácia, sensibilidade,
        AUC, etc.) devem ser transcritas exatamente como publicadas.
  \item \textbf{Imparcialidade:} A análise crítica deve apontar tanto
        contribuições quanto limitações, sem viés de confirmação.
  \item \textbf{SDD:} O mapeamento SDD deve explicitar o problema
        odontológico, os tipos de dado de entrada/saída e os critérios de
        aceitação reportados no artigo.
  \item \textbf{Confiança:} Atribuir valor entre 0 (mínima) e 1 (máxima)
        refletindo a confiabilidade do estudo com base no desenho, tamanho
        amostral, validação e reprodutibilidade.
\end{itemize}

\section*{Verificação de DOIs}
Antes de finalizar o fichamento, verifique se o DOI do artigo está ativo
em \url{https://doi.org/}. DOIs que retornam \texttt{HTTP 404} ou
\texttt{HTTP 403} devem ser sinalizados no campo \textbf{DOI} com a
observação \texttt{(pendente -- possível DOI fabricado)} e reportados no
arquivo \texttt{00-doi-audit.tex}.

\section*{Integração ao Documento Principal}
Para incluir os fichamentos no livro, adicione no preâmbulo:
\begin{verbatim}
% Inclusão dos fichamentos críticos (apêndice documental)
\input{fichamentos/00-instructions}
\input{fichamentos/fichamento-ODT001}
\input{fichamentos/fichamento-ODT002}
% ...
\end{verbatim}
Ou, alternativamente, crie um apêndice específico:
\begin{verbatim}
\begin{apendicesenv}
\partapendices
\chapter{FICHAMENTOS CRÍTICOS}
\input{fichamentos/00-instructions}
\input{fichamentos/fichamento-ODT001}
% ...
\end{apendicesenv}
\end{verbatim}

\section*{Exemplo de Arquivo}
Consulte \texttt{template-fichamento.tex} para o modelo completo e
\texttt{00-doi-audit.tex} para o relatório de auditoria de DOIs.

\vfill
\begin{center}
\rule{0.5\textwidth}{0.4pt}\\
\small{Documento elaborado conforme NBR 14724 e NBR 6029.}\\
\small{Revisão Qualis A1.}
\end{center}

\input{fichamentos/fichamento-ODT001}
\input{fichamentos/fichamento-ODT002}
% ...
\end{verbatim}
Ou, alternativamente, crie um apêndice específico:
\begin{verbatim}
\begin{apendicesenv}
\partapendices
\chapter{FICHAMENTOS CRÍTICOS}
%==============================================================================
% INSTRUÇÕES DE USO — Fichamentos Críticos
% Livro: Odontologia & Inteligência Artificial
% Conforme NBR 6023 (Referências) e NBR 10520 (Citações)
%==============================================================================

\chapter*{Instruções para Elaboração dos Fichamentos Críticos}
\addcontentsline{toc}{chapter}{Instruções — Fichamentos Críticos}

\section*{Propósito}
Os fichamentos críticos constituem o registro sistematizado da fundamentação
científica do livro. Cada artigo citado via \texttt{\textbackslash citacaoativa}
no texto deve possuir um fichamento correspondente, assegurando
rastreabilidade, auditabilidade e reprodutibilidade das fontes.

\section*{Estrutura do Arquivo}
Cada fichamento deve ser salvo como \texttt{fichamento-ODTXXX.tex}, onde
\texttt{XXX} é o identificador sequencial (001, 002, \ldots, 062).
\begin{itemize}
  \item Utilize o arquivo \texttt{template-fichamento.tex} como base.
  \item Substitua todos os campos entre colchetes pelos valores
        correspondentes do artigo fichado.
  \item A chave bibliográfica (\texttt{CHAVE\_BIB}) deve corresponder à
        entrada no arquivo \texttt{.bib} do projeto.
\end{itemize}

\section*{Seções Obrigatórias}
Cada fichamento deve conter obrigatoriamente as seguintes seções:

\begin{enumerate}
  \item \textbf{Identificação} — Dados bibliográficos completos conforme
        NBR 6023.
  \item \textbf{Objetivo} — Síntese do propósito do artigo em 3--5 linhas.
  \item \textbf{Metodologia} — Desenho do estudo, amostra, métricas e
        validação.
  \item \textbf{Resultados Principais} — Valores numéricos relevantes em
        lista itemizada.
  \item \textbf{Limitações} — Pelo menos duas limitações identificadas.
  \item \textbf{Contribuição para o Livro} — Relação explícita entre o
        artigo e o(s) capítulo(s) que ele fundamenta.
  \item \textbf{Análise Crítica} — Validade interna/externa,
        reprodutibilidade, relevância clínica, conflito de interesses e
        viés.
  \item \textbf{Citações Relevantes} — Transcrição de trechos citados no
        livro com indicação de página.
  \item \textbf{Mapeamento SDD} — Relação problema-entrada-saída-critérios
        conforme metodologia SDD do livro.
  \item \textbf{Referência ABNT} — Referência formatada conforme NBR 6023
        dentro de ambiente \texttt{verbatim}.
\end{enumerate}

\section*{Critérios de Qualidade}
\begin{itemize}
  \item \textbf{Rastreabilidade:} Toda afirmação no livro apoiada por
        \texttt{\textbackslash citacaoativa} deve ter fichamento
        correspondente.
  \item \textbf{Precisão:} Métricas numéricas (acurácia, sensibilidade,
        AUC, etc.) devem ser transcritas exatamente como publicadas.
  \item \textbf{Imparcialidade:} A análise crítica deve apontar tanto
        contribuições quanto limitações, sem viés de confirmação.
  \item \textbf{SDD:} O mapeamento SDD deve explicitar o problema
        odontológico, os tipos de dado de entrada/saída e os critérios de
        aceitação reportados no artigo.
  \item \textbf{Confiança:} Atribuir valor entre 0 (mínima) e 1 (máxima)
        refletindo a confiabilidade do estudo com base no desenho, tamanho
        amostral, validação e reprodutibilidade.
\end{itemize}

\section*{Verificação de DOIs}
Antes de finalizar o fichamento, verifique se o DOI do artigo está ativo
em \url{https://doi.org/}. DOIs que retornam \texttt{HTTP 404} ou
\texttt{HTTP 403} devem ser sinalizados no campo \textbf{DOI} com a
observação \texttt{(pendente -- possível DOI fabricado)} e reportados no
arquivo \texttt{00-doi-audit.tex}.

\section*{Integração ao Documento Principal}
Para incluir os fichamentos no livro, adicione no preâmbulo:
\begin{verbatim}
% Inclusão dos fichamentos críticos (apêndice documental)
\input{fichamentos/00-instructions}
\input{fichamentos/fichamento-ODT001}
\input{fichamentos/fichamento-ODT002}
% ...
\end{verbatim}
Ou, alternativamente, crie um apêndice específico:
\begin{verbatim}
\begin{apendicesenv}
\partapendices
\chapter{FICHAMENTOS CRÍTICOS}
\input{fichamentos/00-instructions}
\input{fichamentos/fichamento-ODT001}
% ...
\end{apendicesenv}
\end{verbatim}

\section*{Exemplo de Arquivo}
Consulte \texttt{template-fichamento.tex} para o modelo completo e
\texttt{00-doi-audit.tex} para o relatório de auditoria de DOIs.

\vfill
\begin{center}
\rule{0.5\textwidth}{0.4pt}\\
\small{Documento elaborado conforme NBR 14724 e NBR 6029.}\\
\small{Revisão Qualis A1.}
\end{center}

\input{fichamentos/fichamento-ODT001}
% ...
\end{apendicesenv}
\end{verbatim}

\section*{Exemplo de Arquivo}
Consulte \texttt{template-fichamento.tex} para o modelo completo e
\texttt{00-doi-audit.tex} para o relatório de auditoria de DOIs.

\vfill
\begin{center}
\rule{0.5\textwidth}{0.4pt}\\
\small{Documento elaborado conforme NBR 14724 e NBR 6029.}\\
\small{Revisão Qualis A1.}
\end{center}

\input{fichamentos/fichamento-ODT001}
\input{fichamentos/fichamento-ODT002}
% ...
\end{verbatim}
Ou, alternativamente, crie um apêndice específico:
\begin{verbatim}
\begin{apendicesenv}
\partapendices
\chapter{FICHAMENTOS CRÍTICOS}
%==============================================================================
% INSTRUÇÕES DE USO — Fichamentos Críticos
% Livro: Odontologia & Inteligência Artificial
% Conforme NBR 6023 (Referências) e NBR 10520 (Citações)
%==============================================================================

\chapter*{Instruções para Elaboração dos Fichamentos Críticos}
\addcontentsline{toc}{chapter}{Instruções — Fichamentos Críticos}

\section*{Propósito}
Os fichamentos críticos constituem o registro sistematizado da fundamentação
científica do livro. Cada artigo citado via \texttt{\textbackslash citacaoativa}
no texto deve possuir um fichamento correspondente, assegurando
rastreabilidade, auditabilidade e reprodutibilidade das fontes.

\section*{Estrutura do Arquivo}
Cada fichamento deve ser salvo como \texttt{fichamento-ODTXXX.tex}, onde
\texttt{XXX} é o identificador sequencial (001, 002, \ldots, 062).
\begin{itemize}
  \item Utilize o arquivo \texttt{template-fichamento.tex} como base.
  \item Substitua todos os campos entre colchetes pelos valores
        correspondentes do artigo fichado.
  \item A chave bibliográfica (\texttt{CHAVE\_BIB}) deve corresponder à
        entrada no arquivo \texttt{.bib} do projeto.
\end{itemize}

\section*{Seções Obrigatórias}
Cada fichamento deve conter obrigatoriamente as seguintes seções:

\begin{enumerate}
  \item \textbf{Identificação} — Dados bibliográficos completos conforme
        NBR 6023.
  \item \textbf{Objetivo} — Síntese do propósito do artigo em 3--5 linhas.
  \item \textbf{Metodologia} — Desenho do estudo, amostra, métricas e
        validação.
  \item \textbf{Resultados Principais} — Valores numéricos relevantes em
        lista itemizada.
  \item \textbf{Limitações} — Pelo menos duas limitações identificadas.
  \item \textbf{Contribuição para o Livro} — Relação explícita entre o
        artigo e o(s) capítulo(s) que ele fundamenta.
  \item \textbf{Análise Crítica} — Validade interna/externa,
        reprodutibilidade, relevância clínica, conflito de interesses e
        viés.
  \item \textbf{Citações Relevantes} — Transcrição de trechos citados no
        livro com indicação de página.
  \item \textbf{Mapeamento SDD} — Relação problema-entrada-saída-critérios
        conforme metodologia SDD do livro.
  \item \textbf{Referência ABNT} — Referência formatada conforme NBR 6023
        dentro de ambiente \texttt{verbatim}.
\end{enumerate}

\section*{Critérios de Qualidade}
\begin{itemize}
  \item \textbf{Rastreabilidade:} Toda afirmação no livro apoiada por
        \texttt{\textbackslash citacaoativa} deve ter fichamento
        correspondente.
  \item \textbf{Precisão:} Métricas numéricas (acurácia, sensibilidade,
        AUC, etc.) devem ser transcritas exatamente como publicadas.
  \item \textbf{Imparcialidade:} A análise crítica deve apontar tanto
        contribuições quanto limitações, sem viés de confirmação.
  \item \textbf{SDD:} O mapeamento SDD deve explicitar o problema
        odontológico, os tipos de dado de entrada/saída e os critérios de
        aceitação reportados no artigo.
  \item \textbf{Confiança:} Atribuir valor entre 0 (mínima) e 1 (máxima)
        refletindo a confiabilidade do estudo com base no desenho, tamanho
        amostral, validação e reprodutibilidade.
\end{itemize}

\section*{Verificação de DOIs}
Antes de finalizar o fichamento, verifique se o DOI do artigo está ativo
em \url{https://doi.org/}. DOIs que retornam \texttt{HTTP 404} ou
\texttt{HTTP 403} devem ser sinalizados no campo \textbf{DOI} com a
observação \texttt{(pendente -- possível DOI fabricado)} e reportados no
arquivo \texttt{00-doi-audit.tex}.

\section*{Integração ao Documento Principal}
Para incluir os fichamentos no livro, adicione no preâmbulo:
\begin{verbatim}
% Inclusão dos fichamentos críticos (apêndice documental)
%==============================================================================
% INSTRUÇÕES DE USO — Fichamentos Críticos
% Livro: Odontologia & Inteligência Artificial
% Conforme NBR 6023 (Referências) e NBR 10520 (Citações)
%==============================================================================

\chapter*{Instruções para Elaboração dos Fichamentos Críticos}
\addcontentsline{toc}{chapter}{Instruções — Fichamentos Críticos}

\section*{Propósito}
Os fichamentos críticos constituem o registro sistematizado da fundamentação
científica do livro. Cada artigo citado via \texttt{\textbackslash citacaoativa}
no texto deve possuir um fichamento correspondente, assegurando
rastreabilidade, auditabilidade e reprodutibilidade das fontes.

\section*{Estrutura do Arquivo}
Cada fichamento deve ser salvo como \texttt{fichamento-ODTXXX.tex}, onde
\texttt{XXX} é o identificador sequencial (001, 002, \ldots, 062).
\begin{itemize}
  \item Utilize o arquivo \texttt{template-fichamento.tex} como base.
  \item Substitua todos os campos entre colchetes pelos valores
        correspondentes do artigo fichado.
  \item A chave bibliográfica (\texttt{CHAVE\_BIB}) deve corresponder à
        entrada no arquivo \texttt{.bib} do projeto.
\end{itemize}

\section*{Seções Obrigatórias}
Cada fichamento deve conter obrigatoriamente as seguintes seções:

\begin{enumerate}
  \item \textbf{Identificação} — Dados bibliográficos completos conforme
        NBR 6023.
  \item \textbf{Objetivo} — Síntese do propósito do artigo em 3--5 linhas.
  \item \textbf{Metodologia} — Desenho do estudo, amostra, métricas e
        validação.
  \item \textbf{Resultados Principais} — Valores numéricos relevantes em
        lista itemizada.
  \item \textbf{Limitações} — Pelo menos duas limitações identificadas.
  \item \textbf{Contribuição para o Livro} — Relação explícita entre o
        artigo e o(s) capítulo(s) que ele fundamenta.
  \item \textbf{Análise Crítica} — Validade interna/externa,
        reprodutibilidade, relevância clínica, conflito de interesses e
        viés.
  \item \textbf{Citações Relevantes} — Transcrição de trechos citados no
        livro com indicação de página.
  \item \textbf{Mapeamento SDD} — Relação problema-entrada-saída-critérios
        conforme metodologia SDD do livro.
  \item \textbf{Referência ABNT} — Referência formatada conforme NBR 6023
        dentro de ambiente \texttt{verbatim}.
\end{enumerate}

\section*{Critérios de Qualidade}
\begin{itemize}
  \item \textbf{Rastreabilidade:} Toda afirmação no livro apoiada por
        \texttt{\textbackslash citacaoativa} deve ter fichamento
        correspondente.
  \item \textbf{Precisão:} Métricas numéricas (acurácia, sensibilidade,
        AUC, etc.) devem ser transcritas exatamente como publicadas.
  \item \textbf{Imparcialidade:} A análise crítica deve apontar tanto
        contribuições quanto limitações, sem viés de confirmação.
  \item \textbf{SDD:} O mapeamento SDD deve explicitar o problema
        odontológico, os tipos de dado de entrada/saída e os critérios de
        aceitação reportados no artigo.
  \item \textbf{Confiança:} Atribuir valor entre 0 (mínima) e 1 (máxima)
        refletindo a confiabilidade do estudo com base no desenho, tamanho
        amostral, validação e reprodutibilidade.
\end{itemize}

\section*{Verificação de DOIs}
Antes de finalizar o fichamento, verifique se o DOI do artigo está ativo
em \url{https://doi.org/}. DOIs que retornam \texttt{HTTP 404} ou
\texttt{HTTP 403} devem ser sinalizados no campo \textbf{DOI} com a
observação \texttt{(pendente -- possível DOI fabricado)} e reportados no
arquivo \texttt{00-doi-audit.tex}.

\section*{Integração ao Documento Principal}
Para incluir os fichamentos no livro, adicione no preâmbulo:
\begin{verbatim}
% Inclusão dos fichamentos críticos (apêndice documental)
\input{fichamentos/00-instructions}
\input{fichamentos/fichamento-ODT001}
\input{fichamentos/fichamento-ODT002}
% ...
\end{verbatim}
Ou, alternativamente, crie um apêndice específico:
\begin{verbatim}
\begin{apendicesenv}
\partapendices
\chapter{FICHAMENTOS CRÍTICOS}
\input{fichamentos/00-instructions}
\input{fichamentos/fichamento-ODT001}
% ...
\end{apendicesenv}
\end{verbatim}

\section*{Exemplo de Arquivo}
Consulte \texttt{template-fichamento.tex} para o modelo completo e
\texttt{00-doi-audit.tex} para o relatório de auditoria de DOIs.

\vfill
\begin{center}
\rule{0.5\textwidth}{0.4pt}\\
\small{Documento elaborado conforme NBR 14724 e NBR 6029.}\\
\small{Revisão Qualis A1.}
\end{center}

\input{fichamentos/fichamento-ODT001}
\input{fichamentos/fichamento-ODT002}
% ...
\end{verbatim}
Ou, alternativamente, crie um apêndice específico:
\begin{verbatim}
\begin{apendicesenv}
\partapendices
\chapter{FICHAMENTOS CRÍTICOS}
%==============================================================================
% INSTRUÇÕES DE USO — Fichamentos Críticos
% Livro: Odontologia & Inteligência Artificial
% Conforme NBR 6023 (Referências) e NBR 10520 (Citações)
%==============================================================================

\chapter*{Instruções para Elaboração dos Fichamentos Críticos}
\addcontentsline{toc}{chapter}{Instruções — Fichamentos Críticos}

\section*{Propósito}
Os fichamentos críticos constituem o registro sistematizado da fundamentação
científica do livro. Cada artigo citado via \texttt{\textbackslash citacaoativa}
no texto deve possuir um fichamento correspondente, assegurando
rastreabilidade, auditabilidade e reprodutibilidade das fontes.

\section*{Estrutura do Arquivo}
Cada fichamento deve ser salvo como \texttt{fichamento-ODTXXX.tex}, onde
\texttt{XXX} é o identificador sequencial (001, 002, \ldots, 062).
\begin{itemize}
  \item Utilize o arquivo \texttt{template-fichamento.tex} como base.
  \item Substitua todos os campos entre colchetes pelos valores
        correspondentes do artigo fichado.
  \item A chave bibliográfica (\texttt{CHAVE\_BIB}) deve corresponder à
        entrada no arquivo \texttt{.bib} do projeto.
\end{itemize}

\section*{Seções Obrigatórias}
Cada fichamento deve conter obrigatoriamente as seguintes seções:

\begin{enumerate}
  \item \textbf{Identificação} — Dados bibliográficos completos conforme
        NBR 6023.
  \item \textbf{Objetivo} — Síntese do propósito do artigo em 3--5 linhas.
  \item \textbf{Metodologia} — Desenho do estudo, amostra, métricas e
        validação.
  \item \textbf{Resultados Principais} — Valores numéricos relevantes em
        lista itemizada.
  \item \textbf{Limitações} — Pelo menos duas limitações identificadas.
  \item \textbf{Contribuição para o Livro} — Relação explícita entre o
        artigo e o(s) capítulo(s) que ele fundamenta.
  \item \textbf{Análise Crítica} — Validade interna/externa,
        reprodutibilidade, relevância clínica, conflito de interesses e
        viés.
  \item \textbf{Citações Relevantes} — Transcrição de trechos citados no
        livro com indicação de página.
  \item \textbf{Mapeamento SDD} — Relação problema-entrada-saída-critérios
        conforme metodologia SDD do livro.
  \item \textbf{Referência ABNT} — Referência formatada conforme NBR 6023
        dentro de ambiente \texttt{verbatim}.
\end{enumerate}

\section*{Critérios de Qualidade}
\begin{itemize}
  \item \textbf{Rastreabilidade:} Toda afirmação no livro apoiada por
        \texttt{\textbackslash citacaoativa} deve ter fichamento
        correspondente.
  \item \textbf{Precisão:} Métricas numéricas (acurácia, sensibilidade,
        AUC, etc.) devem ser transcritas exatamente como publicadas.
  \item \textbf{Imparcialidade:} A análise crítica deve apontar tanto
        contribuições quanto limitações, sem viés de confirmação.
  \item \textbf{SDD:} O mapeamento SDD deve explicitar o problema
        odontológico, os tipos de dado de entrada/saída e os critérios de
        aceitação reportados no artigo.
  \item \textbf{Confiança:} Atribuir valor entre 0 (mínima) e 1 (máxima)
        refletindo a confiabilidade do estudo com base no desenho, tamanho
        amostral, validação e reprodutibilidade.
\end{itemize}

\section*{Verificação de DOIs}
Antes de finalizar o fichamento, verifique se o DOI do artigo está ativo
em \url{https://doi.org/}. DOIs que retornam \texttt{HTTP 404} ou
\texttt{HTTP 403} devem ser sinalizados no campo \textbf{DOI} com a
observação \texttt{(pendente -- possível DOI fabricado)} e reportados no
arquivo \texttt{00-doi-audit.tex}.

\section*{Integração ao Documento Principal}
Para incluir os fichamentos no livro, adicione no preâmbulo:
\begin{verbatim}
% Inclusão dos fichamentos críticos (apêndice documental)
\input{fichamentos/00-instructions}
\input{fichamentos/fichamento-ODT001}
\input{fichamentos/fichamento-ODT002}
% ...
\end{verbatim}
Ou, alternativamente, crie um apêndice específico:
\begin{verbatim}
\begin{apendicesenv}
\partapendices
\chapter{FICHAMENTOS CRÍTICOS}
\input{fichamentos/00-instructions}
\input{fichamentos/fichamento-ODT001}
% ...
\end{apendicesenv}
\end{verbatim}

\section*{Exemplo de Arquivo}
Consulte \texttt{template-fichamento.tex} para o modelo completo e
\texttt{00-doi-audit.tex} para o relatório de auditoria de DOIs.

\vfill
\begin{center}
\rule{0.5\textwidth}{0.4pt}\\
\small{Documento elaborado conforme NBR 14724 e NBR 6029.}\\
\small{Revisão Qualis A1.}
\end{center}

\input{fichamentos/fichamento-ODT001}
% ...
\end{apendicesenv}
\end{verbatim}

\section*{Exemplo de Arquivo}
Consulte \texttt{template-fichamento.tex} para o modelo completo e
\texttt{00-doi-audit.tex} para o relatório de auditoria de DOIs.

\vfill
\begin{center}
\rule{0.5\textwidth}{0.4pt}\\
\small{Documento elaborado conforme NBR 14724 e NBR 6029.}\\
\small{Revisão Qualis A1.}
\end{center}

\input{fichamentos/fichamento-ODT001}
% ...
\end{apendicesenv}
\end{verbatim}

\section*{Exemplo de Arquivo}
Consulte \texttt{template-fichamento.tex} para o modelo completo e
\texttt{00-doi-audit.tex} para o relatório de auditoria de DOIs.

\vfill
\begin{center}
\rule{0.5\textwidth}{0.4pt}\\
\small{Documento elaborado conforme NBR 14724 e NBR 6029.}\\
\small{Revisão Qualis A1.}
\end{center}

\input{fichamentos/fichamento-ODT001}
% ...
\end{apendicesenv}
\end{verbatim}

\section*{Exemplo de Arquivo}
Consulte \texttt{template-fichamento.tex} para o modelo completo e
\texttt{00-doi-audit.tex} para o relatório de auditoria de DOIs.

\vfill
\begin{center}
\rule{0.5\textwidth}{0.4pt}\\
\small{Documento elaborado conforme NBR 14724 e NBR 6029.}\\
\small{Revisão Qualis A1.}
\end{center}
